\documentclass[11pt]{article}
\usepackage{amsmath,amssymb,amsthm}
\newtheorem{theorem}{Theorem}
\newtheorem{lemma}{Lemma}
\title{A Leading-Coefficient Audit for Weighted Exterior Sign Cubes}
\author{}
\date{}
\begin{document}
\maketitle

\paragraph{Scope and verification.}
This is a complete written argument, with a finite exact linear-algebra
checker, not a Lean theorem. It concerns one fixed finite pole pool
$\varepsilon x_h$, one fixed analytic weight vector, and fixed query
times that avoid the zeros of the relevant leading coefficients.
It allows the weights to depend analytically on $\varepsilon$.
It does not cover queries at those zeros, moving queries, higher-order
pole deformations, or weights chosen anew for each desired sign word.
In particular it proves neither accepted global HeadComplexity target.

Fix $H\geq1$, real numbers $x_h$, and a real analytic vector
$w(\varepsilon)=\sum_{k\geq0}\varepsilon^k w_k$ near zero. For
$s\in\{-1,1\}^H$ consider
\[
 f_s(t,\varepsilon)=\sum_{h=1}^H
       \frac{s_h w_h(\varepsilon)}{t-\varepsilon x_h}.
 \tag{1}
\]
All query times below are fixed and positive, so all denominators are
positive for sufficiently small positive $\varepsilon$.
Write $T=\operatorname{diag}(x_1,\ldots,x_H)$ and
\[
 D_\ell=\operatorname{span}\{T^j w_k:j+k\leq\ell\},
 \qquad D_{-1}=\{0\},\qquad
 a_\ell=\dim D_\ell-\dim D_{\ell-1}.
 \tag{2}
\]

\begin{lemma}[Slow increase of new leading directions]
One has $a_0\leq1$, $a_\ell\leq a_{\ell-1}+1$ for $\ell\geq1$,
and $\sum_{\ell\geq0}a_\ell\leq H$. Consequently every $a_\ell$ is
at most
\[
 d_H=\left\lfloor\frac{\sqrt{8H+1}-1}{2}\right\rfloor.
 \tag{3}
\]
\end{lemma}
\begin{proof}
The identity
$D_\ell=D_{\ell-1}+TD_{\ell-1}+\operatorname{span}(w_\ell)$
follows directly from (2). Since $TD_{\ell-2}\subseteq D_{\ell-1}$,
the image of $TD_{\ell-1}$ in $D_\ell/D_{\ell-1}$ has dimension at
most $a_{\ell-1}$. The new vector $w_\ell$ adds at most one dimension.
The sum telescopes and is bounded by the ambient dimension $H$.
If $a_\ell=d$, the preceding increments are at least
$d-1,d-2,\ldots,1$. These exist because $a_0\leq1$ and each step
increases by at most one. Therefore $H\geq d(d+1)/2$, giving (3).
\end{proof}

The coefficient of $\varepsilon^\ell$ in (1) is
\[
 t^{-\ell-1} P_{\ell,s}(t),\qquad
 P_{\ell,s}(t)=\sum_{k=0}^{\ell}
                t^k\langle s,T^{\ell-k}w_k\rangle.
 \tag{4}
\]
Thus all coefficient functions of orders below $\ell$ vanish
identically in $t$ exactly when $s\perp D_{\ell-1}$.
On this annihilator, the possible polynomials $P_{\ell,s}$ form a
linear space of dimension at most $a_\ell$: their coefficients depend
only on the restriction of the linear functional $\langle s,-\rangle$
to $D_\ell/D_{\ell-1}$. If the first nonzero coefficient occurs at
order $\ell$, then necessarily $a_\ell>0$.

\begin{theorem}[Generic fixed-query pattern bound]
Fix positive query times $t_1,\ldots,t_q$. Suppose that for each
$s\in\{-1,1\}^H$ with a nonzero coefficient function in (4), its
first nonzero polynomial $P_{\ell,s}$ is nonzero at every $t_j$.
For all sufficiently small positive $\varepsilon$, the number of
strict sign words
\[
 (\operatorname{sign}f_s(t_1,\varepsilon),\ldots,
  \operatorname{sign}f_s(t_q,\varepsilon))
\]
is at most
\[
 \sum_{\ell:a_\ell>0}
       2\sum_{i=0}^{a_\ell-1}\binom{q-1}{i}
 \ \leq\ 2H\sum_{i=0}^{d_H-1}\binom{q-1}{i}.
 \tag{5}
\]
Here binomial coefficients outside their usual range are zero.
\end{theorem}
\begin{proof}
At a fixed first order $\ell$, evaluation of the at most
$a_\ell$-dimensional polynomial space at $q$ points is a linear map.
The strict sign words are therefore bounded by the number of regions
of $q$ central hyperplanes in dimension $a_\ell$. The standard
insertion recurrence, adding one hyperplane and counting the regions
it cuts, gives at most
$2\sum_{i=0}^{a_\ell-1}\binom{q-1}{i}$ regions. Degeneracies only
reduce this bound.

The hypothesis on leading values ensures that each nonzero score
has its leading sign at every query for sufficiently small positive
$\varepsilon$. There are finitely many inputs and queries, so one
common threshold for $\varepsilon$ works. If every coefficient
function is zero, analyticity makes the score identically zero near
zero and it contributes no strict word. Finally at most $H$ increments
are nonzero, each is at most $d_H$, and summing proves (5).
\end{proof}

For every fixed $c>0$, taking $q\geq cH/\log H$ with
$q=O(H/\log H)$ makes the logarithm of the bound in (5)
$O(\sqrt H\log H)=o(q)$. Hence this generic-query regime cannot even
produce $\exp(c' q)$ distinct words for a fixed $c'>0$, for all large
$H$. The assumption about query zeros is substantive. A leading
polynomial can vanish at a prescribed query and expose a later
coefficient there; the proof does not bound that refinement. This is
precisely an available route for a different construction, rather
than a concluded obstruction to it.

\paragraph{Exact checks.}
The companion Python script checks (4), the filtration identities,
and the increment inequalities on rational examples. It also checks
examples with increments $1,2,\ldots,d$ in dimension $d(d+1)/2$,
showing that the square-root bound on an individual increment is
sharp for this general filtration. These finite checks supplement
the written proof and are not used as a universality argument.
\end{document}
